\documentclass[11pt,a4paper,twoside,openright]{memoir}

\usepackage[english]{babel}
\usepackage{fontspec}
\usepackage{microtype}
\usepackage[hidelinks]{hyperref}

\IfFontExistsTF{Libertinus Serif}{%
  \setmainfont{Libertinus Serif}
  \setsansfont{Libertinus Sans}[Scale=0.94]
}{%
  \setmainfont{Latin Modern Roman}
  \setsansfont{Latin Modern Sans}
}

\setstocksize{297mm}{210mm}
\settrimmedsize{\stockheight}{\stockwidth}{*}
\settypeblocksize{220mm}{125mm}{*}
\setlrmargins{30mm}{*}{1.35}
\setulmargins{28mm}{*}{1.2}
\checkandfixthelayout

\chapterstyle{hangnum}
\pagestyle{headings}

\setlength{\parindent}{1.2em}
\setlength{\parskip}{0pt}

\begin{document}

\frontmatter
\title{Classic Editorial}
\author{Author}
\date{}
\maketitle
\tableofcontents

\mainmatter
\chapter{Long-form reading}

A classic editorial system should support sustained reading without demanding attention
for itself. Measure, leading, margins, folios, chapter openings, and page parity work
together as a durable system.

\section{Proportion}
The page should preserve a comfortable text block and meaningful inner and outer margins.
For duplex documents, binding and page parity are functional requirements rather than
cosmetic choices.

\section{Typography}
Use a text face suited to long reading. Keep heading interventions restrained and establish
continuity through consistent paragraph rhythm.

\chapter{Recurring structure}
Repeated chapter architecture should remain stable throughout a book. Variation belongs
to content, not arbitrary styling.

\backmatter
\chapter{Notes}
Back matter should remain subordinate to the primary reading experience while being easy
to navigate.

\end{document}
